% !TEX root = ./Basilisk-avsLibrary20170812.tex

\section{User Guide}

\subsection{{\tt linearAlgebra} Library}
The linear algebra library provides numerous C-based functions to perform basic matrix math operations.  For a complete list of functions supported, consult the {\tt linearAlgebra.h} file.

\subsubsection{Vector Operations}
The vector related functions all begin with a letter '{\tt v}' and are broken down in the following categories depending on if the linear algebra operation is performed on a matrix of a general size, or of a specific size of 2, 3, 4 or 6.
\begin{itemize}
	\item if the function is  {\tt vXXXX()} the code works on a $n$x1 matrix of arbitrary length $n$.
	\item if the function is {\tt v2XXXX()}  the code works on a matrix of length 2 defined as {\tt double vec[2]}.
	\item if the function is {\tt v3XXXX()}  the code works on a matrix of length 2 defined as {\tt double vec[3]}.
	\item if the function is {\tt v4XXXX()} the code works on a matrix of length 2 defined as {\tt double vec[4]}.
\end{itemize}

The most extensive set of linear algebra support functions are for operation on 3 and 4 dimensional matrices common with orbital mechanics (position and velocity vector representations) and attitude dynamics (working with 3- and 4-parameter attitude descriptions).  The following list details what vector operation each function performs:
\begin{itemize}
	\item {\tt Copy}: Copies vector $\bm v_{1}$ into vector $\bm v_{2}$
	\item {\tt SetZero}: Returns a zero vector
	\item {\tt Set}: Sets the vector representation to specific values
	\item {\tt Add}: Sums  $\bm v_{1}$ and  $\bm v_{2}$
	\item {\tt Subtract}: Returns the difference  $\bm v_{1}$ - $\bm v_{2}$
	\item {\tt Scale}: Returns $\bm v_{1}$ multiplied by a scalar $\alpha$
	\item {\tt Dot}: Returns the dot product $\bm v_{1} \cdot \bm v_{2}$
	\item {\tt OuterProduct}: Returns the outer product $\bm v_{1} \bm v_{2}^{T}$
	\item {\tt Normalize}: Returns a normalized vector $\bm v_{1}/ v_{1}$
	\item {\tt MaxAbs}: Returns the largest element of a vector
	\item {\tt IsEqual}: Checks if two vector represnetations are identical
	\item {\tt IsZero}: Checks if a vector is full of zero elements
	\item {\tt Print}: Prints the vector representation to a file
	\item {\tt Sort}: Sorts the vector elements by size
	\item {\tt Print}: Prints the vector representation to a file
	\item {\tt PrintScreen}: Prints the vector representation to the screen
\end{itemize}





\subsubsection{Matrix Operations}
The matrix related functions are all labeled with a letter '{\tt M}' and are broken down into the following matrix dimension related categories:
\begin{itemize}
	\item  {\tt mXXXXy()}: the code works on a $n\times m$ matrix of arbitrary dimension $n$ and $m$.
	\item {\tt m22XXXX()}: the code works on a $2\times 2$ matrix defined through {\tt mat[2][2]}.
	\item {\tt m33XXXX()}: the code works on a $3\times 3$ matrix defined through {\tt mat[3]3]}.
	\item {\tt m44XXXX()}: the code works on a $4\times 4$ matrix defined through {\tt mat[4][4]}.
\end{itemize}

The following list provides an overview of the supported matrix functions.  Note that not al dimensions have all functions provided, but the 2, 3, and 4 dimensional matrix support is pretty complete.
\begin{itemize}
	\item {\tt SetIdentity}: Returns an identity matrix
	\item {\tt SetZero}: Returns a zero matrix
	\item {\tt Set}: Creates a matrix with specific values
	\item {\tt Copy}: Copies the matrix $[M_{1}]$ into $[M_{2}]$
	\item {\tt m33MultM33}: Performs the matrix to matrix multiplication $[M_{1}][M_{2}]$
	\item {\tt m33MultM33t}: Performs the matrix to matrix multiplication $[M_{1}][M_{2}]^{T}$
	\item {\tt m33tMultM33}: Performs the matrix to matrix multiplication $[M_{1}]^{T}[M_{2}]$
	\item {\tt m33MultV3}: Computes $[M_{1}]\bm v_{1}$
	\item {\tt m33tMultV3}: Computes $[M_{1}]^{T}\bm v_{1}$
	\item {\tt v3tMultM33}: Computes $\bm v_{1}^{T}[M_{1}]$
	\item {\tt v3tMultM33t}: Computes  $\bm v_{1}^{T}[M_{1}]^{T}$
	\item {\tt Tilde}: Returns the skew-symmetric matrix $[\tilde{\bm v}_{1}]$
	\item {\tt Transpose}: Returns $[M_{1}]^{T}$
	\item {\tt Add}: Returns the sum $[M_{1}] + [M_{2}]$
	\item {\tt Subtract}: Returns the difference $[M_{1}] - [M_{2}]$
	\item {\tt Scale}: Returns the scaled matrix $\alpha [M_{1}]$
	\item {\tt Trace}: Compute the matrix trace $\sum M_{ii}$
	\item {\tt Determinant}: Returns the square matrix determinant
	\item {\tt IsEqual}: Checks if all matrix elements are equal
	\item {\tt IsZero}: Checks if all matrix elements are zero
	\item {\tt Print}: Prints the matrix representation to a file
	\item {\tt PrintScreen}: Prints the matrix representation to the screen
	\item {\tt Inverse}: Returns the matrix inverse $[M_{1}]^{-1}$
	\item {\tt SingularValues}: Computes the singular values of $[M_{1}]$
	\item {\tt EigenValues}: Computes the Eigenvalues of $[M_{1}]$
	\item {\tt ConditionNumber}: Computes the condition number of $[M_{1}]$
\end{itemize}




\subsection{{\tt rigidBodyKinematics} Library}

The following discussion is a brief overview of the {\tt rigidBodyKinematics} library function notation.  Please see Appendix E in Reference~\citenum{schaub} for a complete description:
\begin{itemize}
	\item {\tt XXX2YYY}: Converts the attitude coordinates {\tt XXX} to {\tt YYY}
	\item {\tt addXXX}: Add the two attitude descriptions $\bm x_{1} = \bm x_{\cal B/F}$ and $\bm x_{2}= \bm x_{\cal F/N}$ to return the sequential rotation of first $\bm x_{1}$ and then $\bm x_{2}$.  Returns is the description $\bm x_{3} = \bm x_{\cal B/N}$
	\item {\tt subXXX}: Subtract the two attitude descriptions $\bm x_{1} = \bm x_{\cal F/N}$ and $\bm x_{2}= \bm x_{\cal B/N}$ to return the relative orientation $\bm x_{3} = \bm x_{\cal F/B}$
	\item {\tt BmatXXX}: Returns the matrix $[B]$ of the differential kinematic equations of the attitude parameters {\tt XXX} in Reference~\citenum{schaub}.  Note the scalar multiplier is not included here.  For example, with MRPs\cite{survey} $[B]$ is defined as $$\dot{\bm \sigma}_{\cal B/N} = \frac{1}{4}[B(\bm\sigma_{\cal B/N})] \leftexp{B}{\bm\omega}_{\cal B/N}$$
	\item {\tt BinvXXX}: Returns the matrix inverse $[B]^{-1}$ of the attitude parameters {\tt XXX}
	\item {\tt dX}: Returns the time derivatives of the attitude parameters {\tt XXX} as a function of these parameters and the body angular velocity vector $\bm \omega$
	\item {\tt Mi}: Returns the single-axis DCM about the $\hat{\bm b}_{i}$ axis
	\item {\tt wrapToPi}: Makes sure that an angle  lies within $\pm\pi$
\end{itemize}




\subsection{{\tt orbitalMotion} Library}

\subsubsection{Orbit Anomaly Angle Conversions}
To convert between the various anomaly angles, the following functions are defined:
\begin{itemize}
	\item {\tt double f2E(double f, double e)}
	\item {\tt double E2f(double Ecc, double e)}
	\item {\tt double E2M(double Ecc, double e)}
	\item {\tt double M2E(double M, double e)}, using bracketed Newton iteration with a maximum of 200 iterations
	\item {\tt double f2H(double f, double e)}
	\item {\tt double H2f(double H, double e)}
	\item {\tt double H2N(double H, double e)}
	\item {\tt double N2H(double N, double e)}, using bracketed Newton iteration with a maximum of 200 iterations
\end{itemize}
Both solvers stop when the residual divided by the equation derivative is within
four times machine epsilon of the anomaly magnitude, or when the root bracket
cannot be refined at floating-point precision. Invalid inputs and iteration
exhaustion produce an error diagnostic and return NaN.



\subsubsection{Orbit Element Conversion}
To convert from classical orbit elements to inertial Cartesian coordinates, the function
\begin{verbatim}
	elem2rv(double mu, ClassicElements *elements, double *rVec, double *vVec)
\end{verbatim}
is used where {\tt mu} is the gravitational constant of the 2-body problem, the classical elements are defined through
$$
	(a, e, i, \Omega, \omega, f)
$$
Inclination must lie in $0\le i\le\pi$, including both equators.
The anomaly is normally the true anomaly $f$. Eccentricity $e=1$ alone does
not identify a parabola: rectilinear states also have unit eccentricity.
A nonrectilinear parabola uses $a=0$, $e=1$, and positive {\tt rPeriap}.
A rectilinear ellipse uses $a>0$, $e=1$, and eccentric anomaly $E$ in the
{\tt f} field; a rectilinear hyperbola uses $a<0$ and hyperbolic anomaly $H$.
A rectilinear parabola uses $a=0$, $e=1$, {\tt rPeriap}=0, positive
{\tt rmag}, and {\tt f}=0 for outward motion or $\pi$ for inward motion.  The function returns the inertial position and velocity vectors in the arrays {\tt rVec} and {\tt vVec}.

To convert from  inertial Cartesian coordinates to classical orbit elements, the function
\begin{verbatim}
	rv2elem(double mu, double *rVec, double *vVec, ClassicElements *elements)
\end{verbatim}
Beyond the classical elements listed above, this routine also stores periapsis
and apoapsis radii $r_p$, $r_a$, and $\alpha=1/a$.
Python retains NumPy scalar fields and its default planar parabolic encoding
$a=-r_p$, {\tt rApoap}=$-1$.
Optionally use {\tt rv2elem(mu, r, v, canonical=True)} for the C-compatible encoding.
The Python {\tt *\_parab} adapters preserve their signed angles and the
inward-positive radial elliptic anomaly convention; negate $E$ when explicitly
converting those inputs for the canonical {\tt elem2rv} convention. Both
conventions remain supported without deprecation warnings or a removal schedule.
The forward adapter interprets every $e=1$, $a<0$ input as a planar parabola
with $r_p=-a$, independently of auxiliary fields. Use the core {\tt elem2rv}
to reconstruct rectilinear hyperbolas. Python and SWIG input elements are not modified.


\paragraph{Mean/Osculating and Equinoctial Domains}
The checked first-order $J_2$ mapping requires finite elliptic inputs
($a>0$, $0\le e<1$, $0\le i\le\pi$), positive equatorial radius, and direction
$+1$ or $-1$. Nonzero $J_2$ rejects
$|1-5\cos^2 i|\le\sqrt{\epsilon}$, where $\epsilon$ is double-precision machine
epsilon. Mapped elements must remain finite and elliptic; invalid results are
not clipped into range. Zero $J_2$ copies the six valid input elements exactly,
including at critical inclinations. These numerical guards do not establish
perturbation accuracy: corrections must still be small. Use
{\tt clMeanOscMapChecked} for these checks.

The prograde equinoctial conversion {\tt clElem2eqElemChecked} also requires finite
elliptic inputs. It supports circular and prograde equatorial orbits but is
singular at the retrograde equator, rejecting $\cos(i/2)\le8\epsilon$.
Coordinates close to that boundary remain ill-conditioned.
Both routines emit diagnostics and all-NaN outputs on failure in C; their
Python counterparts raise {\tt ValueError} without changing the output object.

\paragraph{Hill Frame Domain}
The {\tt hillFrameChecked}, {\tt hill2rvChecked}, and
{\tt rv2hillChecked} routines require finite
three-vectors and a nonzero chief position and velocity with a resolved orbital
plane. They reject
$\|\hat{\bm r}\mathbin{\times}\hat{\bm v}\|\le8\epsilon$.
Normalizing before crossing avoids dependence on distance units. State
transformations additionally require a representable positive angular rate and
finite outputs. Invalid inputs produce all-NaN outputs with a diagnostic in C
or {\tt ValueError} in Python. Distances must use consistent units, velocities
the corresponding units per second.

\paragraph{Compatibility}
The original mapping and Hill function names use the improved calculation in
its supported domain. Rejected domains retain historical formulas, exceptions,
and sentinel outputs. Inclination bounds are always enforced. Original methods
remain supported without deprecation warnings or a removal schedule; checked
calls are optional additions.

\subsubsection{Space Environment Functions}
\paragraph{Atmospheric Density}
The Earth's atmospheric density $\rho$ is computed using\\
\indent	{\tt double atmosphericDensity(double alt)} \\
The density is returned as a scalar value.

\paragraph{Mean Debye Length}
The mean Debye Length for the near-Earth environment is approximated, very crudely, through a polynomial fit.  The function returns the scalar $\lambda_{d}$ and is called through \\
\indent {\tt double debyeLength(double alt)}

\paragraph{Atmospheric Drag Acceleration}
To compute an estimate of  the Earth's atmospheric drag acceleration, use the function: \\
\begin{verbatim}
void atmosphericDrag(double Cd, double A, double m,
                     double *rvec, double *vvec, double *advec);
\end{verbatim}

The inputs are the dimensionless drag coefficient {\tt Cd}, projected area {\tt A}
in m$^2$, and spacecraft mass {\tt m} in kg.  The inertial position {\tt rvec}
is in km, velocity {\tt vvec} in km/s, and returned acceleration {\tt advec} in
km/s$^2$.  The atmosphere is assumed at rest in this inertial frame; no rotation
or wind correction is applied.  The density fit is valid from 100 to 1000 km;
positive altitudes outside that range use the existing extrapolation, including
an exponential tail above 1000 km rather than a hard cutoff.

All inputs and vector norms must be finite, with {\tt Cd} and {\tt A} nonnegative,
{\tt m} positive, and position above Earth's equatorial radius.  Valid zero speed,
area, coefficient, or computed density gives zero acceleration.  Invalid inputs
and acceleration overflow return three NaNs with a diagnostic in C and raise
{\tt ValueError} in Python.  Output underflow is allowed.

\paragraph{Gravitational Zonal Harmonics}
The header safely supplies Earth for the three-argument C call:
\begin{verbatim}
jPerturb(rvec, num, ajtot);
\end{verbatim}
Select another body with the typed interface:
\begin{verbatim}
void jPerturbForPlanet(double *rvec, int num, double *ajtot,
                      CelestialObject_t planetID);
jPerturbForPlanet(rvec, 2, ajtot, CELESTIAL_VENUS);
\end{verbatim}
Four-argument {\tt jPerturb} calls remain supported without deprecation warnings;
the typed function above provides optional checked selection. The bare variadic symbol is
retained for explicit-planet function pointers, which require a fourth integer
argument. Three-argument pointers should target {\tt jPerturbEarth}.
Both source call forms accept C compound literals and C++ braced initializers.
In Python, the zonal utility retains its string planet argument and defaults
to Earth. Set {\tt strict=True} to request checked selection.

Position is in km in a planet-centered frame whose $z$ axis is the spin axis;
acceleration is returned in that frame in km/s$^2$. Earth supports degrees
2--6; other supported bodies provide $J_2$ only. Checked calls return zero for
Sun and Pluto, while Phobos and Deimos are unsupported. The degree must be
2--6 and position must be finite and nonzero. Unsupported bodies, invalid
inputs, and nonfinite acceleration produce a diagnostic and three NaNs in C,
or {\tt ValueError} in Python's checked path. The original
planet selection retains Earth fallback for unknown bodies, unchanged native
output for Sun/Pluto, and Python's {\tt None} Sun/Pluto sentinel and floating
degree cutoffs. These conventions remain supported.

\paragraph{Solar Radiation Pressure Acceleration}
To compute the inertial disturbance acceleration due to the solar radiation pressure use  the function\\
\indent {\tt void solarRad(double A, double m, double *sunvec, double *arvec)}

Here {\tt A} is the projected surface area, {\tt m} is the spacecraft mass, {\tt sunvec} is the sun position vector to the Sun in units of AU.
